\documentclass[12pt,a4paper]{article}

\usepackage{amsmath,amssymb,amsfonts}
\usepackage{geometry}
\usepackage{hyperref}
\usepackage{booktabs}
\usepackage{array}
\usepackage{xcolor}
\usepackage{physics}
\usepackage{graphicx}
\usepackage{enumitem}
\usepackage{fancyhdr}
\usepackage{titlesec}
\usepackage{caption}
\usepackage{listings}
\usepackage{mdframed}

\geometry{margin=2.5cm}

\hypersetup{
    colorlinks=true,
    linkcolor=blue!60!black,
    citecolor=blue!60!black,
    urlcolor=blue!60!black
}

\titleformat{\section}{\large\bfseries}{}{0em}{\thesection\quad}
\titleformat{\subsection}{\normalsize\bfseries}{}{0em}{\thesubsection\quad}

\pagestyle{fancy}
\fancyhf{}
\rhead{\small Exclusive Pion Electroproduction Kinematics}
\lhead{\small \leftmark}
\cfoot{\thepage}

% Box for important results
\newmdenv[
  backgroundcolor=blue!5,
  linecolor=blue!40,
  linewidth=1pt,
  innerleftmargin=10pt,
  innerrightmargin=10pt,
  innertopmargin=8pt,
  innerbottommargin=8pt
]{importantbox}

% Box for code/algorithm
\newmdenv[
  backgroundcolor=gray!8,
  linecolor=gray!50,
  linewidth=1pt,
  innerleftmargin=10pt,
  innerrightmargin=10pt,
  innertopmargin=8pt,
  innerbottommargin=8pt
]{codebox}

\begin{document}

% ============================================================
%  Title
% ============================================================
\begin{center}
  {\LARGE\bfseries Kinematics of Exclusive Pion Electroproduction}\\[0.8em]
  {\large\bfseries Forward Production and the Mandelstam Variable $t$}\\[1.2em]
  {\normalsize Reference document for \texttt{pionCT\_kinematics.C}}\\[0.4em]
  {\small\color{gray} Derived from analysis of JLab pion electroproduction experiments}
\end{center}

\vspace{1em}
\hrule
\vspace{1.5em}

\tableofcontents

\newpage

% ============================================================
\section{Reaction Overview}
% ============================================================

We consider the exclusive charged-pion electroproduction reaction in the
one-photon-exchange approximation:
\begin{equation}
  e(k) \;+\; p(P) \;\longrightarrow\; e'(k') \;+\; \pi^+(p_\pi) \;+\; n(P'),
  \label{eq:reaction}
\end{equation}
where the quantities in parentheses are the four-momenta of each particle.
The incoming electron has energy $E$ and scatters to energy $E'$ at polar
angle $\theta_e$ relative to the beam axis. The target proton is at rest in
the laboratory frame.

\subsection{Four-momenta}

Working in natural units ($c = \hbar = 1$), the relevant four-momenta are:
\begin{align}
  k   &= (E,\, \vec{k})                  && \text{incoming electron,}\\
  k'  &= (E',\, \vec{k}')                && \text{scattered electron,}\\
  q   &= k - k' = (\nu,\, \vec{q}\,)    && \text{virtual photon,}\\
  P   &= (M_p,\, \vec{0}\,)             && \text{target proton (at rest),}\\
  p_\pi &= (E_\pi,\, \vec{p}_\pi)        && \text{produced pion,}\\
  P'  &                                  && \text{recoil neutron.}
\end{align}
The particle masses used throughout are
\begin{equation}
  M_p = 0.938272\;\text{GeV},\qquad
  M_n = 0.939565\;\text{GeV},\qquad
  M_\pi = 0.139571\;\text{GeV}.
\end{equation}

% ============================================================
\section{Lorentz Invariants}
% ============================================================

\subsection{Virtuality \texorpdfstring{$Q^2$}{Q2}}

The four-momentum transfer from the electron to the hadronic system defines
the virtuality of the exchanged photon:
\begin{equation}
  Q^2 \;\equiv\; -q^2 \;=\; 4 E E' \sin^2\!\frac{\theta_e}{2}.
  \label{eq:Q2}
\end{equation}
$Q^2$ is defined as a positive quantity. It controls the resolution of the
virtual photon probe: large $Q^2$ means short-wavelength, hard scattering.

\subsection{Energy transfer and photon three-momentum}

\begin{align}
  \nu &\;\equiv\; E - E' \label{eq:nu}\\[4pt]
  |\vec{q}\,| &\;=\; \sqrt{\nu^2 + Q^2} \label{eq:qmag}
\end{align}

\subsection{Hadronic invariant mass \texorpdfstring{$W$}{W}}

$W$ is the invariant mass of the hadronic final state ($\pi n$ system):
\begin{equation}
  W^2 \;\equiv\; (P + q)^2 \;=\; M_p^2 + 2 M_p \nu - Q^2.
  \label{eq:W2}
\end{equation}
Given $Q^2$, $E$, and a chosen $W$, the scattered electron energy follows
directly:
\begin{equation}
  \nu = \frac{W^2 - M_p^2 + Q^2}{2 M_p},
  \qquad
  E' = E - \nu.
  \label{eq:Eprime}
\end{equation}
Pion production requires $W > M_p + M_\pi$.

\subsection{Electron scattering angle}

Inverting Eq.~\eqref{eq:Q2}:
\begin{equation}
  \theta_e = 2\arcsin\!\sqrt{\frac{Q^2}{4 E E'}}.
  \label{eq:thetae}
\end{equation}

% ============================================================
\section{Virtual Photon Three-Vector in the Lab}
% ============================================================

The beam defines the $+z$ direction. The electron scatters in the $(x,z)$
plane (we choose $\phi = 0$):
\begin{align}
  \vec{k}  &= (0,\; 0,\; E),\\
  \vec{k}' &= (E' \sin\theta_e,\; 0,\; E' \cos\theta_e).
\end{align}
The virtual photon three-vector is therefore
\begin{equation}
  \vec{q} = \vec{k} - \vec{k}' =
  \begin{pmatrix} -E' \sin\theta_e \\ 0 \\ E - E' \cos\theta_e \end{pmatrix}.
  \label{eq:qvec}
\end{equation}
Note that $|\vec{q}\,|$ computed from Eq.~\eqref{eq:qvec} is numerically
identical to $\sqrt{\nu^2 + Q^2}$ (Eq.~\eqref{eq:qmag}), providing a useful
cross-check. The unit vector along $\vec{q}$ and its in-plane perpendicular
are:
\begin{equation}
  \hat{q} = \frac{\vec{q}}{|\vec{q}\,|},
  \qquad
  \hat{q}_\perp = \begin{pmatrix} -q_z/|\vec{q}\,| \\ 0 \\ q_x/|\vec{q}\,| \end{pmatrix}.
  \label{eq:qhat}
\end{equation}

% ============================================================
\section{Two-Body CM Kinematics}
% ============================================================

The hadronic system is treated as a two-body final state:
\begin{equation}
  \gamma^*(q) + p(P) \;\longrightarrow\; \pi(p_\pi) + n(P').
\end{equation}

\subsection{CM pion momentum}

In the hadronic CM frame the pion momentum magnitude $p^*$ is fixed by $W$
alone (standard two-body kinematics):
\begin{equation}
  \boxed{
  p^* \;\equiv\; |\vec{p}_\pi^{\,*}|
     \;=\; \frac{1}{2W}
       \sqrt{\bigl[W^2 - (M_p + M_\pi)^2\bigr]\,\bigl[W^2 - (M_p - M_\pi)^2\bigr]}.
  }
  \label{eq:pstar}
\end{equation}
The corresponding CM pion energy is $E_\pi^* = \sqrt{(p^*)^2 + M_\pi^2}$.

\subsection{CM pion four-vector}

With the CM polar angle $\theta^*$ (measured from the $\vec{q}$ direction)
and azimuthal angle $\phi^* = 0$, the pion four-vector in the CM is:
\begin{equation}
  p_\pi^{*\,\mu} = \bigl(E_\pi^*,\; p^* \sin\theta^*,\; 0,\; p^* \cos\theta^*\bigr),
  \label{eq:piCM}
\end{equation}
where the $z$-axis in the CM is aligned with $\vec{q}$.

% ============================================================
\section{Lorentz Boost to the Laboratory}
% ============================================================

\subsection{Boost velocity}

The hadronic CM frame moves along $\vec{q}$ with velocity
\begin{equation}
  \beta_\text{CM} = \frac{|\vec{q}\,|}{M_p + \nu},
  \qquad
  \gamma_\text{CM} = \frac{1}{\sqrt{1 - \beta_\text{CM}^2}}.
  \label{eq:betaCM}
\end{equation}
This is a \emph{vector} boost along $\hat{q}$, not along the beam axis.

\subsection{Decomposition in the CM}

The pion momentum in the CM is decomposed along and perpendicular to $\hat{q}$:
\begin{align}
  p_\parallel^* &= p^* \cos\theta^*  && \text{(along } \hat{q}\text{)},\\
  p_\perp^*     &= p^* \sin\theta^*  && \text{(perpendicular to } \hat{q}\text{)}.
\end{align}

\subsection{Boost equations}

The Lorentz boost along $\hat{q}$ transforms the components as:
\begin{align}
  p_\parallel^\text{lab} &= \gamma_\text{CM}\bigl(p_\parallel^* + \beta_\text{CM}\,E_\pi^*\bigr),
  \label{eq:boost_par}\\
  E_\pi^\text{lab}       &= \gamma_\text{CM}\bigl(E_\pi^* + \beta_\text{CM}\,p_\parallel^*\bigr),
  \label{eq:boost_E}\\
  p_\perp^\text{lab}     &= p_\perp^*.
  \label{eq:boost_perp}
\end{align}
The perpendicular component is unchanged by the boost.

\subsection{Lab Cartesian components}

Projecting back onto the beam frame $(x,z)$ axes using $\hat{q}$ and
$\hat{q}_\perp$ from Eq.~\eqref{eq:qhat}:
\begin{align}
  p_{\pi,x}^\text{lab} &= p_\parallel^\text{lab}\,\hat{q}_x
                          + p_\perp^\text{lab}\,\hat{q}_{\perp,x},
  \label{eq:ppix}\\
  p_{\pi,z}^\text{lab} &= p_\parallel^\text{lab}\,\hat{q}_z
                          + p_\perp^\text{lab}\,\hat{q}_{\perp,z}.
  \label{eq:ppiz}
\end{align}

\begin{importantbox}
\textbf{Important:} This vector-boost recipe is exact and requires no
subsequent rotation. A common error is to boost along $+z$ first and then
apply a \texttt{RotateY} correction --- this is only an approximation and
introduces small errors when $\theta_e$ is large.
\end{importantbox}

% ============================================================
\section{Laboratory Pion Observables}
% ============================================================

\subsection{Momentum and angle}

\begin{align}
  |\vec{p}_\pi| &= \sqrt{p_{\pi,x}^2 + p_{\pi,z}^2},
  \label{eq:pmag}\\[4pt]
  \theta_\pi &= \arctan\!\left(\frac{|p_{\pi,x}|}{p_{\pi,z}}\right),
  \label{eq:thetapi}
\end{align}
where $\theta_\pi$ is the polar angle of the pion relative to the beam axis.

\subsection{Longitudinal momentum surplus \texorpdfstring{$k_\pi$}{kpi}}

The quantity $k_\pi$ listed in kinematic tables is:
\begin{equation}
  k_\pi \;\equiv\; |\vec{q}\,| - |\vec{p}_\pi|.
  \label{eq:kpi}
\end{equation}
At $\theta^* = 0$ (forward production), the pion travels exactly along
$\vec{q}$, so $k_\pi = |\vec{q}\,| - |\vec{p}_\pi|$ equals the longitudinal
momentum surplus of the virtual photon over the pion. It is approximately
$\sqrt{|t_\text{min}|}$ for small $|t|$.

% ============================================================
\section{The Mandelstam Variable \texorpdfstring{$t$}{t}}
% ============================================================

\subsection{Definition}

\begin{equation}
  t \;\equiv\; (q - p_\pi)^2
    \;=\; (\nu - E_\pi)^2 - |\vec{q} - \vec{p}_\pi|^2.
  \label{eq:t}
\end{equation}
Expanding:
\begin{equation}
  t = (\nu - E_\pi)^2 - (q_x - p_{\pi,x})^2 - (q_z - p_{\pi,z})^2.
  \label{eq:t_expanded}
\end{equation}
$t$ is always negative for physical forward production. It measures the
four-momentum transferred from the virtual photon to the pion, or
equivalently, from the initial proton to the recoil neutron:
$t = (P - P')^2$.

\subsection{Physical interpretation}

\begin{itemize}[leftmargin=2em]
  \item \textbf{$t = 0$:} would require perfect collinearity of photon and
    pion four-momenta; kinematically inaccessible due to mass differences.
  \item \textbf{$t \to t_\text{min}$:} pion goes forward in the CM
    ($\theta^* = 0$), minimum momentum transfer. The cross section peaks here,
    driven by pion-pole exchange.
  \item \textbf{$t \ll t_\text{min}$} (large $|t|$): pion is produced at
    large CM angle; cross section falls steeply, suppressed by form factors.
\end{itemize}

In the near-forward limit, $t$ is related to the CM transverse momentum by:
\begin{equation}
  -t \;\approx\; (k_T^\text{CM})^2 + \text{(small longitudinal term)}
     \;=\; (p^* \sin\theta^*)^2 + \mathcal{O}\!\left(\frac{1}{W^2}\right).
  \label{eq:t_approx}
\end{equation}

\subsection{Minimum value \texorpdfstring{$t_\text{min}$}{tmin}}

$t_\text{min}$ is the value of $t$ at $\theta^* = 0$. It is a hard
kinematic boundary: \emph{every} event at a given $(W, Q^2, E)$ satisfies
$t \leq t_\text{min} < 0$. No event can have $|t| < |t_\text{min}|$.

\begin{importantbox}
$t_\text{min}$ is computed by evaluating Eq.~\eqref{eq:t_expanded} at
$\theta^* = 0$, which gives a pion traveling exactly along $\hat{q}$ in the
lab. It depends on $W$, $Q^2$, and $E$.
\end{importantbox}

Table~\ref{tab:tmin} shows $t_\text{min}$ for the four kinematic settings
at $E = 10.7\;\text{GeV}$.

\begin{table}[h!]
\centering
\caption{Kinematic settings with $E = 10.7\;\text{GeV}$, forward pion
production ($\theta^* = 0$). The $t_\text{min}$ column is the physically
computed value; the $t_\text{table}$ column is what appears in the reference
table. For Q$^2$ = 8.5 the two values differ significantly: $W = 2.95\;\text{GeV}$
and $t = -0.40\;\text{GeV}^2$ cannot coexist.}
\label{tab:tmin}
\begin{tabular}{ccccccccc}
\toprule
$Q^2$ & $W$ & $E'$ & $\theta_e$ & $\theta_\pi$ & $p_\pi$ & $k_\pi$
& $t_\text{min}$ & $t_\text{table}$ \\
(GeV$^2$) & (GeV) & (GeV) & (deg) & (deg) & (GeV/$c$) & (GeV)
& (GeV$^2$) & (GeV$^2$) \\
\midrule
5.0 & 2.42 & 5.384 & 16.94 & 15.79 & 5.103 & 0.665 & $-0.397$ & $-0.40$ \\
6.5 & 2.72 & 3.763 & 23.18 & 11.56 & 6.722 & 0.669 & $-0.402$ & $-0.40$ \\
7.5 & 2.89 & 2.722 & 29.40 &  9.11 & 7.759 & 0.677 & $-0.410$ & $-0.40$ \\
8.5 & 2.95 & 2.002 & 36.72 &  7.50 & 8.451 & 0.723 & $-0.462$ & $-0.40$ \\
\bottomrule
\end{tabular}
\end{table}

% ============================================================
\section{The \texorpdfstring{$t$}{t} Cut in Practice}
% ============================================================

In an exclusive electroproduction analysis, events are accepted within a
range of $t$:
\begin{equation}
  t_\text{min}(W,Q^2,E) \;\leq\; t \;\leq\; t_\text{cut}.
  \label{eq:trange}
\end{equation}
The upper bound $t_\text{cut}$ (closest to zero) is chosen to isolate the
forward-production region where the cross section is dominated by pion-pole
exchange and backgrounds are minimised.

\textbf{When $|t_\text{cut}| < |t_\text{min}|$} (i.e.\ $t_\text{cut} >
t_\text{min}$), the cut window is empty: no events exist with $|t|$ smaller
than $|t_\text{min}|$. This is the situation for $Q^2 = 8.5\;\text{GeV}^2$,
$W = 2.95\;\text{GeV}$ in Table~\ref{tab:tmin}, where $|t_\text{min}| =
0.462 > 0.40$.

% ============================================================
\section{Algorithm Summary}
% ============================================================

The code \texttt{pionCT\_kinematics.C} implements two usage modes.

\subsection{Mode A: $W$ as input}

Given $(Q^2,\, E,\, W)$:
\begin{enumerate}[label=\arabic*.]
  \item Compute $\nu$ and $E'$ from Eqs.~\eqref{eq:W2}--\eqref{eq:Eprime}.
  \item Compute $\theta_e$ from Eq.~\eqref{eq:thetae}.
  \item Build $\vec{q}$ from Eq.~\eqref{eq:qvec}; compute $|\vec{q}\,|$,
    $\hat{q}$, $\hat{q}_\perp$.
  \item Compute $\beta_\text{CM}$, $\gamma_\text{CM}$ from
    Eq.~\eqref{eq:betaCM}.
  \item Compute $p^*$ and $E_\pi^*$ from Eq.~\eqref{eq:pstar}.
  \item At $\theta^* = 0$: set $p_\parallel^* = p^*$, $p_\perp^* = 0$.
  \item Apply boost Eqs.~\eqref{eq:boost_par}--\eqref{eq:boost_perp}.
  \item Compute lab components Eqs.~\eqref{eq:ppix}--\eqref{eq:ppiz}.
  \item Compute $|\vec{p}_\pi|$, $\theta_\pi$, $k_\pi$, $t$ from
    Eqs.~\eqref{eq:pmag}--\eqref{eq:t_expanded}.
  \item Output is $t_\text{min}$ for this $W$.
\end{enumerate}

\subsection{Mode B: $t_\text{min}$ as input}

Given $(Q^2,\, E,\, t_\text{target})$, find $W$ such that
$t_\text{min}(W) = t_\text{target}$:
\begin{enumerate}[label=\arabic*.]
  \item Define $f(W) = t_\text{min}(W) - t_\text{target}$ using Mode A.
  \item Scan $W$ from threshold upward to bracket the zero of $f$.
  \item Apply bisection to find $W$ such that $|f(W)| < 10^{-9}$.
  \item Run Mode A with the found $W$ to produce all observables.
\end{enumerate}

\begin{codebox}
\textbf{Usage examples (ROOT):}\\[4pt]
\texttt{// Mode A: specify W directly}\\
\texttt{pionCT\_kinematics(5.0, 10.7, "W", 2.42)}\\[4pt]
\texttt{// Mode B: find W from t\_min = -0.40}\\
\texttt{pionCT\_kinematics(5.0, 10.7, "t", -0.40)}\\[4pt]
\texttt{// Print full comparison table}\\
\texttt{pionCT\_kinematics()}
\end{codebox}

% ============================================================
\section{Corrected Reference Table}
% ============================================================

Table~\ref{tab:corrected} shows the kinematic settings computed using Mode B
(exact $t_\text{min} = -0.40\;\text{GeV}^2$ for all rows). Note that for
$Q^2 = 8.5\;\text{GeV}^2$ the correct $W$ is $3.091\;\text{GeV}$, not
$2.95\;\text{GeV}$.

\begin{table}[h!]
\centering
\caption{Corrected kinematic settings for $E = 10.7\;\text{GeV}$,
$t_\text{min} = -0.40\;\text{GeV}^2$ exactly. $W$ is the primary output
of the Mode B calculation. The $k_\pi$ value is identical for all rows
by construction.}
\label{tab:corrected}
\begin{tabular}{cccccccc}
\toprule
$Q^2$ & $W$ & $E'$ & $\theta_e$ & $\theta_\pi$ & $p_\pi$ & $k_\pi$ & $t_\text{min}$ \\
(GeV$^2$) & (GeV) & (GeV) & (deg) & (deg) & (GeV/$c$) & (GeV) & (GeV$^2$) \\
\midrule
5.0 & 2.414 & 5.399 & 16.917 & 15.849 & 5.086 & 0.667 & $-0.400$ \\
6.5 & 2.725 & 3.749 & 23.223 & 11.517 & 6.737 & 0.667 & $-0.400$ \\
7.5 & 2.914 & 2.649 & 29.810 &  8.906 & 7.837 & 0.667 & $-0.400$ \\
8.5 & 3.091 & 1.548 & 41.976 &  6.188 & 8.938 & 0.667 & $-0.400$ \\
\bottomrule
\end{tabular}
\end{table}

% ============================================================
\section{Quick Reference: Formula Chain}
% ============================================================

\begin{align}
  Q^2     &= 4EE'\sin^2\!\tfrac{\theta_e}{2}
            \tag{\ref{eq:Q2}}\\
  \nu     &= E - E'
            \tag{\ref{eq:nu}}\\
  |\vec{q}\,| &= \sqrt{\nu^2 + Q^2}
            \tag{\ref{eq:qmag}}\\
  W^2     &= M_p^2 + 2M_p\nu - Q^2
            \tag{\ref{eq:W2}}\\
  p^*     &= \frac{1}{2W}\sqrt{[W^2-(M_p+M_\pi)^2][W^2-(M_p-M_\pi)^2]}
            \tag{\ref{eq:pstar}}\\
  \vec{q} &= (-E'\sin\theta_e,\;0,\;E - E'\cos\theta_e)
            \tag{\ref{eq:qvec}}\\
  \beta_\text{CM} &= \frac{|\vec{q}\,|}{M_p + \nu}
            \tag{\ref{eq:betaCM}}\\
  p_\parallel^\text{lab} &= \gamma(p_\parallel^* + \beta E_\pi^*)
            \tag{\ref{eq:boost_par}}\\
  E_\pi^\text{lab} &= \gamma(E_\pi^* + \beta p_\parallel^*)
            \tag{\ref{eq:boost_E}}\\
  t       &= (\nu - E_\pi)^2 - |\vec{q} - \vec{p}_\pi|^2
            \tag{\ref{eq:t}}\\
  k_\pi   &= |\vec{q}\,| - |\vec{p}_\pi|
            \tag{\ref{eq:kpi}}
\end{align}

\end{document}
